\documentclass[a4paper, 12pt]{article}

\usepackage[utf8]{inputenc}
\usepackage[T1]{fontenc}
\usepackage{textcomp}
\usepackage{amssymb}
\usepackage{newtxtext} \usepackage{newtxmath}
\usepackage{amsmath, amssymb}
\newtheorem{problem}{Problem}
\newtheorem{example}{Example}
\newtheorem{lemma}{Lemma}
\newtheorem{theorem}{Theorem}
\newtheorem{problem}{Problem}
\newtheorem{example}{Example} \newtheorem{definition}{Definition}
\newtheorem{lemma}{Lemma}
\newtheorem{theorem}{Theorem}
\DeclareMathAlphabet{\mathcal}{OMS}{cmsy}{m}{n}
\usepackage{parskip}

\begin{document}

It is with great enthusiasm that I write to strongly support the promotion of
Dr. Goldschmied to the rank of Associate Professor of Psychiatry. As a computer
science student based in Argentina, I have had the privilege of being mentored by
Dr. Goldschmied for the past four years. Even though our collaboration has been
mostly remote, her dedication and guidance have made an impact nothing short of
extraordinary in my academic trajectory. I am honored to provide my highest
recommendation for her promotion, as I know she represents the best of academic
mentorship.

Dr. Goldschmied's mentoring style is defined by her remarkable commitment to
fostering independent academic growth. This was apparent even at the very
beginning of our collaboration. During my second week under her guidance, I came
up with a novel way of quantifying cortical excitability with transcranial
magnetic stimulation to predict depression. I was excited but
nervous to share this with her. After all, it was unrelated to the tasks she had
assigned me, and why would she listen to a junior, remote student?
Many mentors would have redirected me back to routine tasks without
paying much attention, but not Dr. Goldschmied. She responded with enthusiasm,
validation, and genuine interest. She encouraged me to pursue the inquiry,
engaging directly with my idea and providing rigorous and constructive
criticism. Whereas my background in computer science led me to consider it
only in computational or abstract terms, Dr. Goldschmied challenged and trained
me to broaden my perspective and consider its clinical implications as well. She
actively put me in contact with researchers whose specialties were relevant to
the matter, and dedicated significant time to helping me refine the methodology
and draft a manuscript. When it came the time to publish the findings,
recognizing how daunting the submission process can be for a first timer, she
patiently guided me through every step of the way. She even identified and
helped me secure the \textit{DPN Editors’ New Investigator Waiver Award}, which
covered the publication costs for the \textit{DPN Nature} journal. The story I
tell here in a single paragraph comprises three and a half years of dedicated
and patient mentorship, from validating my idea at its earliest stage to helping
me secure the resources to finally see it published. And, all the way through,
Dr. Goldschmied's commitment to championing my journey was unrelenting.


Needless to say, mentoring a student located in Argentina presents particular
challenges, especially when it comes to facilitating connections. Networking is
a crucial aspect of academic life, and one particularly difficult to develop when working
from abroad. Despite this, Dr. Goldschmied consistently found ways to integrate
me into the scientific community. Without any prompting on my part, recognizing
that my interests aligned with their profiles, she connected me with Hans van
Dongen at Washington State University and Diego R. Mazzotti  at the University
of Kansas, with both of whom I developed highly enriching professional
relations. These introductions directly resulted in my visiting Dr.
van Dongen's laboratory in May 2024 and co-authoring a paper with Dr. Mazzotti.
But to complement these direct, virtual connections with large-scale, in-person
exposure, she strongly encouraged and helped me to submit an abstract to the SLEEP 2024
conference. When it was accepted for a poster presentation, she helped me apply
for a grant that ultimately funded my attendance. While at the conference, she
actively introduced me to established researchers whose work aligned with my own
interests. Her efforts to bridge the gap between a student in a developing
nation and the academic community of the US highlights her remarkable commitment
to equity and the success of her mentees.

The most compelling evidence of Dr. Goldschmied’s excellence is the
transformative impact she has had on my career. I entered her lab as an
undergraduate with no practical research experience, no scientific network, and
only a rudimentary understanding of neuroscience. She patiently guided my
education from the ground up, training me in techniques such as EEG analysis and
TMS, as well as grounding me in theoretical concepts such as intracortical
inhibition/facilitation, slow-wave activity, and the sleep homeostasis
hypothesis. It is not an overstatement to characterize her impact on my academic
trajectory as radical. I have since published my first first-author paper
(with Dr. Goldschmied as senior author), presented a poster at the SLEEP 2024
conference, and am preparing an oral presentation for SLEEP 2026, where my
attendance will be funded by Dr. Goldschmied's lab. Thanks to her tireless
efforts, I am now well-positioned to apply for doctoral programs and turn what
once seemed like a distant dream into a reality.

Dr. Goldschmied took an inexperienced, international student from the Global
South with nothing but raw ambition, and meticulously developed me into a
capable, published researcher. Without a doubt, she embodies the very best of
academic mentorship. That is why I offer my highest and most enthusiastic
recommendation for her promotion to Associate Professor of Psychiatry.

Sincerely,

Santiago López Pereyra,\newline
Department of Mathematics, Astronomy, Physics, and Computer Science\newline
National University of Córdoba, Argentina\newline
santiago.lopez.pereyra@mi.unc.edu.ar




\end{document}



